\documentclass[a4,11pt]{aleph-notas}
% Para obtener solo el pdf, compilar con pdfLaTeX.
% Para obtener el xml compilar con XeLaTeX.

% -- Paquetes adicionales
% \usepackage[handout]{aleph-moodle}
\usepackage{aleph-moodle}
\moodleregisternewcommands
% Todos los comandos nuevos deben ir luego del comando anterior
\usepackage{aleph-comandos}


% -- Datos
\institucion{Facultad de Ciencias Exactas, Naturales y Ambientales}
\carrera{Catálogo STEM}
\asignatura{Matemática III}
\tema{Cuestionario No. 2: Derivada en varias variables}
\autor{Andrés Merino}
\fecha{Periodo 2026-1}

\logouno[0.16\textwidth]{Logos/logoPUCE_04_ac}
\definecolor{colortext}{HTML}{1957d1}
\definecolor{colordef}{HTML}{1957d1}
\fuente{montserrat}


\begin{document}

\encabezado

\vspace*{-8mm}
\tableofcontents

%%%%%%%%%%%%%%%%%%%%%%%%%%%%%%%%%%%%%%%%
\section{Indicaciones}
%%%%%%%%%%%%%%%%%%%%%%%%%%%%%%%%%%%%%%%%

Se plantean bancos de preguntas orientados a evaluar el \textbf{criterio}: «Interpreta el concepto de derivada en varias variables y reconocer sus elementos y propiedades», correspondiente al \textbf{resultado de aprendizaje}: Identificar los conceptos, teoremas y propiedades de funciones vectoriales, derivadas parciales e integral vectorial.

%%%%%%%%%%%%%%%%%%%%%%%%%%%%%%%%%%%%%%%%
\section{Banco de preguntas}
%%%%%%%%%%%%%%%%%%%%%%%%%%%%%%%%%%%%%%%%


%%%%%%%%%%%%%%%%%%%%%%%%%%%%%%%%%%%%%%%%
\begin{quiz}{Derivadas parciales: concepto e interpretación}
%%%%%%%%%%%%%%%%%%%%%%%%%%%%%%%%%%%%%%%%

\begin{multi}[]%
    {DP-Concepto-01}
    Sea $f$ un campo escalar. ¿Qué tipo de objeto es $\dfrac{\partial f}{\partial x}$?
    \item* Un campo escalar
    \item Un campo vectorial
    \item Una matriz
    \item Un número real
\end{multi}

\begin{multi}[]%
    {DP-Concepto-01b}
    Sea $f:\R^2\to\R$ un campo escalar diferenciable. ¿Qué tipo de objeto es $\nabla f$?
    \item* Un campo vectorial
    \item Un campo escalar
    \item Una matriz de $2\times 2$
    \item Un número real
\end{multi}

\begin{multi}[]%
    {DP-Concepto-02}
    La derivada parcial $\dfrac{\partial f}{\partial x}(a,b)$ representa geométricamente:
    \item* La pendiente en la dirección de la primera variable en $(a,b)$.
    \item La pendiente del plano tangente a la gráfica de $f$ en $(a,b)$.
    \item La tasa de cambio de $f$ en la dirección del vector gradiente.
    \item La longitud del vector normal a la superficie en $(a,b)$.
\end{multi}

\begin{multi}[]%
    {DP-Concepto-02b}
    La derivada parcial $\dfrac{\partial f}{\partial y}(a,b)$ representa geométricamente:
    \item* La pendiente en la dirección de la segunda variable en $(a,b)$.
    \item La pendiente del plano tangente a la gráfica de $f$ en $(a,b)$.
    \item La tasa de cambio de $f$ en la dirección del gradiente.
    \item El valor máximo de $f$ en $(a,b)$.
\end{multi}

\begin{multi}[]%
    {DP-Concepto-03}
    Para calcular $\dfrac{\partial f}{\partial x}(x,y)$ se considera $y$ como:
    \item* Una constante y se deriva respecto de $x$.
    \item Una variable y se aplica la regla del producto.
    \item La variable independiente principal.
    \item Una función de $x$.
\end{multi}

\begin{multi}[]%
    {DP-Concepto-03b}
    Para calcular $\dfrac{\partial f}{\partial y}(x,y)$ se considera $x$ como:
    \item* Una constante y se deriva respecto de $y$.
    \item Una variable y se aplica la regla del cociente.
    \item La variable independiente principal.
    \item Una función de $y$.
\end{multi}

\begin{multi}[]%
    {DP-Concepto-04}
    El teorema de Schwarz establece que, bajo condiciones de regularidad:
    \item* $\dfrac{\partial^2 f}{\partial x \partial y} = \dfrac{\partial^2 f}{\partial y \partial x}$
    \item $\dfrac{\partial f}{\partial x} = \dfrac{\partial f}{\partial y}$
    \item $\dfrac{\partial^2 f}{\partial x^2} = \dfrac{\partial^2 f}{\partial y^2}$
    \item $\dfrac{\partial^2 f}{\partial x \partial y} = \dfrac{\partial f}{\partial x} \cdot \dfrac{\partial f}{\partial y}$
\end{multi}

\begin{multi}[]%
    {DP-Concepto-04b}
    Si $f$ tiene derivadas parciales de segundo orden continuas, el teorema de Schwarz implica que la matriz hessiana es:
    \item* Simétrica
    \item Diagonal
    \item Antisimétrica
    \item Invertible
\end{multi}

\begin{multi}[]%
    {DP-Concepto-05}
    Si $f$ es un campo escalar diferenciable en $(a,b)$, ¿cuál de las siguientes afirmaciones es correcta?
    \item* $f$ es continua en $(a,b)$.
    \item $f$ tiene todas sus derivadas parciales de orden superior en $(a,b)$.
    \item Las derivadas parciales de $f$ son continuas en $(a,b)$.
    \item $f$ es diferenciable en todo $\R^2$.
\end{multi}

\begin{multi}[]%
    {DP-Concepto-05b}
    Si $f$ tiene derivadas parciales en $(a,b)$ pero no es diferenciable en $(a,b)$, entonces:
    \item* $f$ puede ser discontinua en $(a,b)$.
    \item $f$ es necesariamente continua en $(a,b)$.
    \item $f$ posee gradiente en $(a,b)$.
    \item $f$ es diferenciable en todo punto cercano a $(a,b)$.
\end{multi}

\begin{multi}[]%
    {DP-Concepto-06}
    La notación $f_x(a,b)$ es equivalente a:
    \item* $\dfrac{\partial f}{\partial x}(a,b)$
    \item $\dfrac{df}{dx}(a,b)$
    \item $f'(a,b)$
    \item $\dfrac{\partial f}{\partial y}(a,b)$
\end{multi}

\begin{multi}[]%
    {DP-Concepto-06b}
    La notación $f_{yy}(a,b)$ es equivalente a:
    \item* $\dfrac{\partial^2 f}{\partial y^2}(a,b)$
    \item $\dfrac{\partial f}{\partial y}(a,b)$
    \item $\left(\dfrac{\partial f}{\partial y}(a,b)\right)^2$
    \item $\dfrac{\partial^2 f}{\partial x \partial y}(a,b)$
\end{multi}

\end{quiz}


%%%%%%%%%%%%%%%%%%%%%%%%%%%%%%%%%%%%%%%%
\begin{quiz}{Cálculo de derivadas parciales: forma simbólica}
%%%%%%%%%%%%%%%%%%%%%%%%%%%%%%%%%%%%%%%%

\begin{multi}[]%
    {DPS-Calc-01}
    Sea $f(x,y) = x^3 y + 2y^2$. Entonces $\dfrac{\partial f}{\partial x}$ es:
    \item* $3x^2 y$
    \item $3x^2 y + 2y^2$
    \item $x^3 + 4y$
    \item $3x^2$
\end{multi}

\begin{multi}[]%
    {DPS-Calc-01c}
    Sea $f(x,y) = x^2 y^3 - 5x$. Entonces $\dfrac{\partial f}{\partial x}$ es:
    \item* $2xy^3 - 5$
    \item $3x^2y^2$
    \item $2xy^3$
    \item $2y^3 - 5$
\end{multi}

\begin{multi}[]%
    {DPS-Calc-01b}
    Sea $f(x,y) = x^3 y + 2y^2$. Entonces $\dfrac{\partial f}{\partial y}$ es:
    \item* $x^3 + 4y$
    \item $3x^2 y$
    \item $x^3 y + 4y$
    \item $3x^2 + 4y$
\end{multi}

\begin{multi}[]%
    {DPS-Calc-01d}
    Sea $f(x,y) = x^2 y^3 - 5x$. Entonces $\dfrac{\partial f}{\partial y}$ es:
    \item* $3x^2 y^2$
    \item $2xy^3$
    \item $3y^2$
    \item $x^2 y^2$
\end{multi}

\begin{multi}[]%
    {DPS-Calc-02}
    Sea $f(x,y) = e^{xy}$. Entonces $\dfrac{\partial f}{\partial x}$ es:
    \item* $y\, e^{xy}$
    \item $e^{xy}$
    \item $x\, e^{xy}$
    \item $xy\, e^{xy-1}$
\end{multi}

\begin{multi}[]%
    {DPS-Calc-02c}
    Sea $f(x,y) = \cos(xy)$. Entonces $\dfrac{\partial f}{\partial x}$ es:
    \item* $-y\,\sin(xy)$
    \item $-\sin(xy)$
    \item $y\,\sin(xy)$
    \item $-xy\,\sin(xy)$
\end{multi}

\begin{multi}[]%
    {DPS-Calc-02b}
    Sea $f(x,y) = \sin(x^2 + y)$. Entonces $\dfrac{\partial f}{\partial y}$ es:
    \item* $\cos(x^2 + y)$
    \item $2x\cos(x^2+y)$
    \item $-\cos(x^2+y)$
    \item $\sin(x^2+y)$
\end{multi}

\begin{multi}[]%
    {DPS-Calc-02d}
    Sea $f(x,y) = \ln(x + y^2)$. Entonces $\dfrac{\partial f}{\partial x}$ es:
    \item* $\dfrac{1}{x+y^2}$
    \item $\dfrac{1}{x}$
    \item $\dfrac{2y}{x+y^2}$
    \item $\ln(1+y^2)$
\end{multi}

\begin{multi}[]%
    {DPS-Calc-03}
    Sea $f(x,y,z) = x^2 y z^3$. Entonces $\dfrac{\partial f}{\partial z}$ es:
    \item* $3x^2 y z^2$
    \item $x^2 y z^2$
    \item $3x^2 z^3$
    \item $3y z^2$
\end{multi}

\begin{multi}[]%
    {DPS-Calc-03b}
    Sea $f(x,y,z) = x\,\sin(yz)$. Entonces $\dfrac{\partial f}{\partial z}$ es:
    \item* $xy\cos(yz)$
    \item $y\cos(yz)$
    \item $x\cos(yz)$
    \item $xy\,\sin(yz)$
\end{multi}

\begin{multi}[]%
    {DPS-Calc-04}
    Sea $f(x,y) = x\ln(y) + y^2$. Entonces $\dfrac{\partial^2 f}{\partial x \partial y}$ es:
    \item* $\dfrac{1}{y}$
    \item $\ln(y)$
    \item $\dfrac{x}{y}$
    \item $\dfrac{x}{y^2}$
\end{multi}

\begin{multi}[]%
    {DPS-Calc-04b}
    Sea $f(x,y) = x^2 e^y$. Entonces $\dfrac{\partial^2 f}{\partial y \partial x}$ es:
    \item* $2x\,e^y$
    \item $x^2 e^y$
    \item $2e^y$
    \item $2x$
\end{multi}

\end{quiz}


%%%%%%%%%%%%%%%%%%%%%%%%%%%%%%%%%%%%%%%%
\begin{quiz}{Cálculo de derivadas parciales: evaluación numérica}
%%%%%%%%%%%%%%%%%%%%%%%%%%%%%%%%%%%%%%%%

\begin{numerical}[tolerance=0.01]%
    {DPN-Calc-01}
    Sea $f(x,y) = x^2 + 3xy$. El valor de $\dfrac{\partial f}{\partial x}(1,2)$ es:
    \item 8
\end{numerical}

\begin{numerical}[tolerance=0.01]%
    {DPN-Calc-05}
    Sea $f(x,y) = x^3 - xy^2$. El valor de $\dfrac{\partial f}{\partial x}(2,1)$ es:
    \item 11
\end{numerical}

\begin{numerical}[tolerance=0.01]%
    {DPN-Calc-01b}
    Sea $f(x,y) = x^2 + 3xy$. El valor de $\dfrac{\partial f}{\partial y}(1,2)$ es:
    \item 3
\end{numerical}

\begin{numerical}[tolerance=0.01]%
    {DPN-Calc-05b}
    Sea $f(x,y) = x^3 - xy^2$. El valor de $\dfrac{\partial f}{\partial y}(2,1)$ es:
    \item -4
\end{numerical}

\begin{numerical}[tolerance=0.01]%
    {DPN-Calc-02}
    Sea $f(x,y) = e^x \sin(y)$. El valor de $\dfrac{\partial f}{\partial y}(0,0)$ es:
    \item 1
\end{numerical}

\begin{numerical}[tolerance=0.01]%
    {DPN-Calc-06}
    Sea $f(x,y) = \sin(x)\cos(y)$. El valor de $\dfrac{\partial f}{\partial x}(0,0)$ es:
    \item 1
\end{numerical}

\begin{numerical}[tolerance=0.01]%
    {DPN-Calc-02b}
    Sea $f(x,y) = e^x \cos(y)$. El valor de $\dfrac{\partial f}{\partial x}(0,0)$ es:
    \item 1
\end{numerical}

\begin{numerical}[tolerance=0.01]%
    {DPN-Calc-06b}
    Sea $f(x,y) = \sin(x)\cos(y)$. El valor de $\dfrac{\partial f}{\partial y}(0,0)$ es:
    \item 0
\end{numerical}

\begin{numerical}[tolerance=0.01]%
    {DPN-Calc-03}
    Sea $f(x,y,z) = x^2 y + yz^2$. El valor de $\dfrac{\partial f}{\partial z}(1,2,3)$ es:
    \item 12
\end{numerical}

\begin{numerical}[tolerance=0.01]%
    {DPN-Calc-07b}
    Sea $f(x,y,z) = xyz^2$. El valor de $\dfrac{\partial f}{\partial x}(1,2,3)$ es:
    \item 18
\end{numerical}

\begin{numerical}[tolerance=0.01]%
    {DPN-Calc-03b}
    Sea $f(x,y,z) = x^2 y + yz^2$. El valor de $\dfrac{\partial f}{\partial y}(2,1,1)$ es:
    \item 5
\end{numerical}

\begin{numerical}[tolerance=0.01]%
    {DPN-Calc-07}
    Sea $f(x,y) = x^2 y - y^3$. El valor de $\dfrac{\partial^2 f}{\partial x^2}(2,1)$ es:
    \item 2
\end{numerical}

\begin{numerical}[tolerance=0.01]%
    {DPN-Calc-04}
    Sea $f(x,y) = \ln(x^2 + y^2)$. El valor de $\dfrac{\partial f}{\partial x}(1,0)$ es:
    \item 2
\end{numerical}

% \begin{numerical}[tolerance=0.01]%
%     {DPN-Calc-08}
%     Sea $f(x,y) = \arctan\!\left(\dfrac{y}{x}\right)$. El valor de $\dfrac{\partial f}{\partial x}(1,1)$ es:
%     \item -0.5
% \end{numerical}

\begin{numerical}[tolerance=0.01]%
    {DPN-Calc-04b}
    Sea $f(x,y) = x^3 - 2x^2 y + y^3$. El valor de $\dfrac{\partial^2 f}{\partial x^2}(1,1)$ es:
    \item 2
\end{numerical}

\begin{numerical}[tolerance=0.01]%
    {DPN-Calc-08b}
    Sea $f(x,y) = x^2 + 2xy + y^2$. El valor de $\dfrac{\partial^2 f}{\partial x \partial y}(3,2)$ es:
    \item 2
\end{numerical}

\end{quiz}


%%%%%%%%%%%%%%%%%%%%%%%%%%%%%%%%%%%%%%%%
\begin{quiz}{Derivada direccional y gradiente}
%%%%%%%%%%%%%%%%%%%%%%%%%%%%%%%%%%%%%%%%

\begin{multi}[]%
    {DG-Dir-01}
    La derivada direccional $f'({x};{u})$ representa:
    \item* La tasa de cambio de $f$ en el punto ${x}$ en la dirección del vector  ${u}$.
    \item El gradiente de $f(x)$ en la dirección ${u}$.
    \item La derivada parcial de $f$ respecto a la primera variable.
    \item La longitud del gradiente de $f$ en ${x}$.
\end{multi}

\begin{multi}[]%
    {DG-Dir-03}
    La derivada direccional máxima de $f$ en un punto ${x}_0$ se alcanza en la dirección de:
    \item* $\nabla f({x}_0)$
    \item $-\nabla f({x}_0)$
    \item Cualquier vector unitario
    \item El vector nulo
\end{multi}

\begin{multi}[]%
    {DG-Dir-02}
    Si $f$ es un campo escalar diferenciable, la relación entre la derivada direccional y el gradiente es:
    \item* $f'({x};{u}) = \nabla f({x}) \cdot {u}$
    \item $f'({x};{u}) = \nabla f({x}) \times {u}$
    \item $f'({x};{u}) = \|\nabla f({x})\| \cdot \|{u}\|$
    \item $f'({x};{u}) = \nabla f({u}) \cdot {x}$
\end{multi}

\begin{multi}[]%
    {DG-Dir-05}
    Si $\|\nabla f({x}_0)\| = 5$ y ${u}$ es un vector unitario en la dirección del gradiente, entonces $f'({x}_0;{u})$ es:
    \item* $5$
    \item $0$
    \item $-5$
    \item $25$
\end{multi}

\begin{multi}[]%
    {DG-Grad-01}
    Si $f:\R^n\to\R$ es un campo escalar, entonces $\nabla f$ es:
    \item* Un campo vectorial
    \item Un campo escalar
    \item Una matriz de $n\times n$
    \item Un número real
\end{multi}

\begin{multi}[]%
    {DG-Grad-03}
    El gradiente $\nabla f({x}_0)$ apunta en dirección perpendicular a:
    \item* Las curvas de nivel de $f$ que pasan por ${x}_0$.
    \item Las tangentes de las trayectorias en ${x}_0$.
    \item El eje $x$.
    \item El vector posición ${x}_0$.
\end{multi}

\begin{multi}[]%
    {DG-Grad-02}
    El gradiente $\nabla f({x}_0)$ apunta en la dirección de:
    \item* Máximo crecimiento de $f$ en ${x}_0$.
    \item Mínimo crecimiento de $f$ en ${x}_0$.
    \item Crecimiento nulo de $f$.
    \item La derivada parcial más grande.
\end{multi}

\begin{multi}[]%
    {DG-Grad-04}
    Si $f(x,y) = x^2 + y^2$, entonces $\nabla f(1,1)$ es:
    \item* $(2,\, 2)$
    \item $(1,\, 1)$
    \item $(2,\, 0)$
    \item $(0,\, 2)$
\end{multi}

\begin{multi}[]%
    {DG-Dir-04}
    Si ${u}$ es perpendicular a $\nabla f({x}_0)$, entonces $f'({x}_0;{u})$ es:
    \item* $0$
    \item $\|\nabla f({x}_0)\|$
    \item $-\|\nabla f({x}_0)\|$
    \item Indefinida
\end{multi}

\begin{multi}[]%
    {DG-Dir-06}
    La derivada direccional mínima de $f$ en ${x}_0$ se obtiene en la dirección de:
    \item* $-\nabla f({x}_0)$
    \item $\nabla f({x}_0)$
    \item Cualquier vector perpendicular al gradiente
    \item El vector posición ${x}_0$
\end{multi}

\end{quiz}


%%%%%%%%%%%%%%%%%%%%%%%%%%%%%%%%%%%%%%%%
\begin{quiz}{Jacobiano y derivada de campos vectoriales}
%%%%%%%%%%%%%%%%%%%%%%%%%%%%%%%%%%%%%%%%

\begin{multi}[]%
    {JC-Def-01}
    Sea $\func{F}{\R^n}{\R^m}$ un campo vectorial diferenciable. La matriz jacobiana $J_F({x})$ es:
    \item* Una matriz de $m\times n$
    \item Una matriz de $n\times m$
    \item Un vector de $\R^m$
    \item Un campo escalar
\end{multi}

\begin{multi}[]%
    {JC-Def-03}
    Sea $\func{F}{\R^2}{\R^3}$. La matriz jacobiana $J_F({x})$ tiene dimensión:
    \item* $3\times 2$
    \item $2\times 3$
    \item $3\times 3$
    \item $2\times 2$
\end{multi}

\begin{multi}[]%
    {JC-Def-02}
    Si $\func{F}{\R^3}{\R^2}$ con $F=(f_1, f_2)$, la entrada $(i,j)$ de la matriz jacobiana es:
    \item* $\dfrac{\partial f_i}{\partial x_j}$
    \item $\dfrac{\partial f_j}{\partial x_i}$
    \item $f_i \cdot x_j$
    \item $\dfrac{\partial^2 f_i}{\partial x_j^2}$
\end{multi}

\begin{multi}[]%
    {JC-Def-04}
    Si $F = (f_1, f_2, f_3)$ con $\func{F}{\R^2}{\R^3}$, la primera columna de $J_F$ contiene:
    \item* Las derivadas parciales de cada $f_i$ respecto a $x_1$
    \item Las derivadas parciales de $f_1$ respecto a $x_1$ y $x_2$
    \item Los valores de $(f_1, f_2, f_3)$ en un punto
    \item Las derivadas parciales de $f_1$ respecto a $x_1, x_2, x_3$
\end{multi}

\begin{multi}[]%
    {JC-Tipo-01}
    Si $\func{F}{\R^n}{\R^m}$ es un campo vectorial y ${a}\in\R^n$, entonces $\dfrac{\partial F}{\partial x_j}({a})$ es:
    \item Un número real
    \item Un campo escalar
    \item* Un campo vectorial
    \item Una matriz
\end{multi}

\begin{multi}[]%
    {JC-Tipo-03}
    Si $\func{F}{\R^n}{\R^m}$, ¿cuántas derivadas parciales tiene cada componente $f_i$?
    \item* $n$
    \item $m$
    \item $nm$
    \item $1$
\end{multi}

\begin{multi}[]%
    {JC-Tipo-02}
    El jacobiano $J_F$ de un campo vectorial $\func{F}{\R^n}{\R^m}$ es:
    \item* Una función de $\R^n$ en el espacio de matrices de $m\times n$
    \item Un campo escalar
    \item Un campo vectorial de $\R^n$ en $\R^m$
    \item Un número real
\end{multi}

\begin{multi}[]%
    {JC-Tipo-04}
    El rango de la matriz jacobiana $J_F({a})$ de $\func{F}{\R^n}{\R^m}$ es como máximo:
    \item* $\min(m,n)$
    \item $m+n$
    \item $mn$
    \item $\max(m,n)$
\end{multi}

\begin{multi}[]%
    {JC-Interp-01}
    La matriz jacobiana $J_F({x}_0)$ representa la mejor aproximación lineal de $F$ cerca de ${x}_0$ en el sentido de que:
    \item* $F({x}_0 + {h}) \approx F({x}_0) + J_F({x}_0)\,{h}$ para ${h}$ pequeño.
    \item $F({x}_0 + {h}) \approx J_F({x}_0)\,{x}_0$ para todo ${h}$.
    \item $F({x}_0 + {h}) = F({x}_0) + J_F({x}_0)\,{h}$ exactamente.
    \item $J_F({x}_0)$ minimiza el error cuadrático de la aproximación.
\end{multi}

\begin{multi}[]%
    {JC-Interp-02}
    Si $J_F({x}_0)$ es invertible, entonces $F$ es localmente:
    \item* Invertible cerca de ${x}_0$ (por el teorema de la función inversa).
    \item Constante cerca de ${x}_0$.
    \item Lineal en todo $\R^n$.
    \item No diferenciable en ${x}_0$.
\end{multi}

\begin{multi}[]%
    {JC-EjCalc-01}
    Sea $F(x,y) = (x^2 + y,\ xy)$. La entrada $(1,2)$ de la matriz jacobiana $JF(x,y)$ es:
    \item* $x$
    \item $y$
    \item $2x$
    \item $1$
\end{multi}

\begin{multi}[]%
    {JC-EjCalc-02}
    Sea $F(x,y) = (x^2 - y^2,\ 2xy)$. La entrada $(2,1)$ de la matriz jacobiana $J_F(x,y)$ es:
    \item* $2y$
    \item $2x$
    \item $-2y$
    \item $2xy$
\end{multi}

\end{quiz}


%%%%%%%%%%%%%%%%%%%%%%%%%%%%%%%%%%%%%%%%
\begin{quiz}{Matriz Hessiana}
%%%%%%%%%%%%%%%%%%%%%%%%%%%%%%%%%%%%%%%%

\begin{multi}[]%
    {H-Def-01}
    Si $f:\R^n\to\R$ es un campo escalar, la matriz hessiana $H_f({x})$ es:
    \item* Una matriz de $n\times n$
    \item Un vector de $\R^n$
    \item Un campo escalar
    \item Una matriz de $1\times n$
\end{multi}

\begin{multi}[]%
    {H-Def-03}
    Si $f:\R^2\to\R$, la diagonal de $H_f$ contiene:
    \item* Las derivadas parciales de segundo orden puras $\dfrac{\partial^2 f}{\partial x_i^2}$
    \item Las derivadas parciales de primer orden
    \item Los valores propios de $H_f$
    \item Las derivadas cruzadas
\end{multi}

\begin{multi}[]%
    {H-Def-02}
    La entrada $(i,j)$ de la matriz hessiana $H_f$ está dada por:
    \item* $\dfrac{\partial^2 f}{\partial x_i \partial x_j}$
    \item $\dfrac{\partial f}{\partial x_i} \cdot \dfrac{\partial f}{\partial x_j}$
    \item $\dfrac{\partial^2 f}{\partial x_i^2}$ solo si $i=j$, y $0$ si $i\neq j$
    \item $\dfrac{\partial f}{\partial x_i} + \dfrac{\partial f}{\partial x_j}$
\end{multi}

\begin{multi}[]%
    {H-Def-04}
    Si $f:\R^3\to\R$, ¿cuántas entradas independientes tiene la matriz hessiana $H_f$ (asumiendo continuidad de las derivadas de segundo orden)?
    \item* $6$
    \item $9$
    \item $3$
    \item $27$
\end{multi}

\begin{multi}[]%
    {H-Prop-01}
    Bajo condiciones de regularidad (continuidad de las derivadas de segundo orden), la matriz hessiana $H_f$ es:
    \item* Simétrica, es decir, $H_f = H_f^\top$
    \item Antisimétrica, es decir, $H_f = -H_f^\top$
    \item Diagonal
    \item Ortogonal
\end{multi}

\begin{multi}[]%
    {H-Prop-02}
    La matriz hessiana $H_f({x}_0)$ es definida positiva si:
    \item* Todos sus valores propios son positivos.
    \item Su traza es positiva.
    \item Su determinante es positivo.
    \item Es simétrica.
\end{multi}

\begin{multi}[]%
    {H-Tipo-01}
    Si $f:\R^n\to\R$, entonces $H_f$ evaluada en un punto ${a}\in\R^n$ es:
    \item* Una matriz de $n\times n$
    \item Un campo escalar
    \item Un campo vectorial
    \item Un vector de $\R^n$
\end{multi}

\begin{multi}[]%
    {H-Tipo-02}
    Si $f:\R^2\to\R$, el determinante de $H_f({a})$ se usa para:
    \item* Clasificar puntos críticos de $f$ como máximos, mínimos o puntos silla.
    \item Calcular el gradiente de $f$ en ${a}$.
    \item Determinar si $f$ es continua en ${a}$.
    \item Hallar los valores propios del jacobiano.
\end{multi}

\end{quiz}


%%%%%%%%%%%%%%%%%%%%%%%%%%%%%%%%%%%%%%%%
\begin{quiz}{Cálculo Hessiana}
%%%%%%%%%%%%%%%%%%%%%%%%%%%%%%%%%%%%%%%%

\begin{multi}[]%
    {H-Calc-01}
    Sea $f(x,y) = x^2 + xy + y^2$. La entrada $(1,2)$ de $H_f$ es:
    \item* $1$
    \item $2$
    \item $0$
    \item $x+y$
\end{multi}

\begin{multi}[]%
    {H-Calc-03}
    Sea $f(x,y) = x^3 - 3xy^2$. La entrada $(2,2)$ de $H_f(x,y)$ es:
    \item* $-6x$
    \item $6x$
    \item $-6y$
    \item $6y$
\end{multi}

\begin{multi}[]%
    {H-Calc-02}
    Sea $f(x,y) = x^2 + xy + y^2$. La traza de $H_f$ (suma de entradas diagonales) es:
    \item* $4$
    \item $2$
    \item $1$
    \item $6$
\end{multi}

\begin{multi}[]%
    {H-Calc-04}
    Sea $f(x,y) = x^2 y^2$. La entrada $(1,2)$ de $H_f(x,y)$ es:
    \item* $4xy$
    \item $2y^2$
    \item $2x^2$
    \item $4x^2y^2$
\end{multi}

\begin{multi}[]%
    {H-Calc-05}
    Sea $f(x,y) = e^{x+y}$. El determinante de $H_f(x,y)$ es:
    \item* $0$
    \item $e^{2(x+y)}$
    \item $2e^{x+y}$
    \item $1$
\end{multi}

\begin{multi}[]%
    {H-Calc-06}
    Sea $f(x,y) = x^2 + y^2$. La traza de $H_f$ es:
    \item* $4$
    \item $2$
    \item $0$
    \item $2x + 2y$
\end{multi}

\begin{multi}[]%
    {H-Calc-07}
    Sea $f(x,y) = \sin(x)\cos(y)$. La entrada $(1,1)$ de $H_f(x,y)$ es:
    \item* $-\sin(x)\cos(y)$
    \item $\cos(x)\cos(y)$
    \item $-\cos(x)\sin(y)$
    \item $\sin(x)\sin(y)$
\end{multi}

\begin{multi}[]%
    {H-Calc-08}
    Sea $f(x,y,z) = x^2 + y^2 + z^2$. La matriz hessiana $H_f$ es:
    \item* $2I_3$ (dos veces la matriz identidad $3\times 3$)
    \item La matriz nula
    \item Una matriz diagonal con entradas $(x,\, y,\, z)$
    \item Una matriz antisimétrica
\end{multi}

\begin{multi}[]%
    {H-Calc-09}
    Sea $f(x,y) = x^2 - 4xy + y^2$. El determinante de $H_f(x,y)$ es:
    \item* $-12$
    \item $4$
    \item $12$
    \item $0$
\end{multi}

\begin{multi}[]%
    {H-Calc-10}
    Sea $f(x,y) = x^4 + y^4$. La entrada $(2,2)$ de $H_f(x,y)$ es:
    \item* $12y^2$
    \item $4y^3$
    \item $12y$
    \item $4y^2$
\end{multi}

\end{quiz}


%%%%%%%%%%%%%%%%%%%%%%%%%%%%%%%%%%%%%%%%
\begin{quiz}{Rotacional y divergencia}
%%%%%%%%%%%%%%%%%%%%%%%%%%%%%%%%%%%%%%%%

\begin{multi}[]%
    {RD-Div-01}
    Si $\func{F}{\R^3}{\R^3}$ es un campo vectorial, entonces $\nabla \cdot F$ es:
    \item* Un campo escalar
    \item Un campo vectorial
    \item Una matriz
    \item Un número real
\end{multi}

\begin{multi}[]%
    {RD-Div-02}
    Si $\func{F}{\R^3}{\R^3}$ con $F=(P,Q,R)$, la fórmula de la divergencia es:
    \item* $\dfrac{\partial P}{\partial x} + \dfrac{\partial Q}{\partial y} + \dfrac{\partial R}{\partial z}$
    \item $\left(\dfrac{\partial R}{\partial y} - \dfrac{\partial Q}{\partial z},\ \dfrac{\partial P}{\partial z} - \dfrac{\partial R}{\partial x},\ \dfrac{\partial Q}{\partial x} - \dfrac{\partial P}{\partial y}\right)$
    \item $P + Q + R$
    \item $\dfrac{\partial P}{\partial x} \cdot \dfrac{\partial Q}{\partial y} \cdot \dfrac{\partial R}{\partial z}$
\end{multi}

\begin{multi}[]%
    {RD-Rot-01}
    Si $\func{F}{\R^3}{\R^3}$ es un campo vectorial, entonces $\nabla \times F$ es:
    \item* Un campo vectorial
    \item Un campo escalar
    \item Una matriz de $3\times 3$
    \item Un número real
\end{multi}

\begin{multi}[]%
    {RD-Rot-02}
    Si $F$ es un campo conservativo (gradiente de algún campo escalar), entonces $\nabla \times F$ es:
    \item* El vector cero $(0,0,0)$
    \item El gradiente de $F$
    \item Un campo escalar no nulo
    \item La divergencia de $F$
\end{multi}

\begin{multi}[]%
    {RD-Dom-01}
    ¿Sobre qué tipo de función está definido el rotacional $\nabla \times F$?
    \item* Campos vectoriales $\func{F}{\R^3}{\R^3}$
    \item Campos escalares $\func{f}{\R^3}{\R}$
    \item Trayectorias $\func{\alpha}{\R}{\R^3}$
    \item Campos vectoriales $\func{F}{\R^2}{\R^2}$
\end{multi}

\begin{multi}[]%
    {RD-Dom-03}
    La divergencia $\nabla \cdot F$ está definida para campos vectoriales $\func{F}{\R^n}{\R^n}$. ¿Qué condición deben cumplir el dominio y el codominio?
    \item* Deben tener la misma dimensión.
    \item El dominio debe ser necesariamente $\R^3$.
    \item El codominio debe ser $\R$.
    \item No hay restricción de dimensión.
\end{multi}

\begin{multi}[]%
    {RD-Dom-02}
    ¿Sobre qué tipo de función está definida la divergencia $\nabla \cdot F$?
    \item* Campos vectoriales $\func{F}{\R^n}{\R^n}$
    \item Campos escalares $\func{f}{\R^n}{\R}$
    \item Trayectorias $\func{\alpha}{\R}{\R^n}$
    \item Campos vectoriales $\func{F}{\R^n}{\R^m}$ con $m\neq n$
\end{multi}

\begin{multi}[]%
    {RD-Dom-04}
    ¿Para cuál de las siguientes funciones está bien definido el rotacional $\nabla \times F$?
    \item* $\func{F}{\R^3}{\R^3}$
    \item $\func{f}{\R^3}{\R}$
    \item $\func{F}{\R^2}{\R^3}$
    \item $\func{\alpha}{\R}{\R^3}$
\end{multi}

\end{quiz}


%%%%%%%%%%%%%%%%%%%%%%%%%%%%%%%%%%%%%%%%
\begin{quiz}{Calculo Rotacional y divergencia}
%%%%%%%%%%%%%%%%%%%%%%%%%%%%%%%%%%%%%%%%

\begin{multi}[]%
    {RD-Calc-01}
    Sea $F(x,y,z) = (x^2, y^2, z^2)$. La divergencia $\nabla \cdot F$ es:
    \item* $2x + 2y + 2z$
    \item $(2x, 2y, 2z)$
    \item $2x \cdot 2y \cdot 2z$
    \item $6$
\end{multi}

\begin{multi}[]%
    {RD-Calc-03}
    Sea $F(x,y,z) = (x, y, z)$. La divergencia $\nabla \cdot F$ es:
    \item* $3$
    \item $0$
    \item $(1,1,1)$
    \item $x+y+z$
\end{multi}

\begin{multi}[]%
    {RD-Calc-02}
    Sea $F(x,y,z) = (yz, xz, xy)$. El rotacional $\nabla \times F$ es:
    \item* $(0, 0, 0)$
    \item $(yz, xz, xy)$
    \item $(1, 1, 1)$
    \item $(x+y, y+z, x+z)$
\end{multi}

\begin{multi}[]%
    {RD-Calc-05}
    Sea $F(x,y,z) = (x^2, y^2, z^2)$. El rotacional $\nabla \times F$ es:
    \item* $(0, 0, 0)$
    \item $(2x, 2y, 2z)$
    \item $(x^2, y^2, z^2)$
    \item $(1, 1, 1)$
\end{multi}

\begin{multi}[]%
    {RD-Calc-04}
    Sea $F(x,y,z) = (xy, yz, zx)$. La divergencia $\nabla \cdot F$ es:
    \item* $x + y + z$
    \item $(y, z, x)$
    \item $0$
    \item $3xyz$
\end{multi}

\begin{multi}[]%
    {RD-Calc-06}
    Sea $F(x,y,z) = (y, -x, 0)$. La divergencia $\nabla \cdot F$ es:
    \item* $0$
    \item $-2$
    \item $2$
    \item $(0, 0, -2)$
\end{multi}

\begin{multi}[]%
    {RD-Calc-07}
    Sea $F(x,y,z) = (y, -x, 0)$. El rotacional $\nabla \times F$ es:
    \item* $(0, 0, -2)$
    \item $(0, 0, 2)$
    \item $(1, -1, 0)$
    \item $(0, 0, 0)$
\end{multi}

\begin{multi}[]%
    {RD-Calc-08}
    Sea $F(x,y,z) = (e^x, e^y, e^z)$. La divergencia $\nabla \cdot F$ es:
    \item* $e^x + e^y + e^z$
    \item $e^{x+y+z}$
    \item $(e^x, e^y, e^z)$
    \item $3\,e^{x+y+z}$
\end{multi}

\begin{multi}[]%
    {RD-Calc-09}
    Sea $F(x,y,z) = (xz, yz, xy)$. La divergencia $\nabla \cdot F$ es:
    \item* $2z$
    \item $xy + yz + xz$
    \item $0$
    \item $(z, z, 0)$
\end{multi}

\begin{multi}[]%
    {RD-Calc-10}
    Si $f(x,y,z) = x^2 + y^2 + z^2$, el laplaciano $\nabla \cdot (\nabla f)$ es:
    \item* $6$
    \item $0$
    \item $(2,2,2)$
    \item $2(x^2+y^2+z^2)$
\end{multi}

\end{quiz}

\end{document}
